\section*{Аннотация}
\addcontentsline{toc}{section}{Аннотация}

Подготовка к единому государственному экзамену по математике требует учебных
материалов, адаптированных под индивидуальный уровень каждого ученика. Существующие
платформы либо предлагают адаптивный подбор задач без генерации объяснительного
контента, либо работают в диалоговом формате без структурированных визуальных
материалов. Ни одна из них не закрывает полный цикл: от аналитики знаний ученика к
персонализированному конспекту с математическими иллюстрациями.

В данной работе разработан модуль генерации таких конспектов с применением больших языковых моделей. На вход поступает
профиль компетенций ученика, построенный по истории решённых задач и допущенных
ошибок; на выходе --- структурированный конспект с текстом и визуальными материалами,
адаптированными под уровень конкретного школьника. Отдельный пайплайн отвечает за
автоматическое создание математических иллюстраций: графиков функций, геометрических
чертежей и блок-схем, --- обеспечивая их геометрическую корректность.

\medskip
\noindent\textbf{Ключевые слова:} персонализированное обучение, генерация учебных
материалов, большие языковые модели, генерация с извлечением
информации

\section*{Abstract}

Preparing for the Unified State Exam in mathematics requires learning materials
tailored to each student's individual knowledge level. Existing platforms either offer
adaptive task selection without generating explanatory content, or operate in a
dialogue format without structured visual materials. None of them covers the full
cycle from student knowledge analytics to a personalized study note with mathematical
illustrations.

This work presents a module that generates such study notes using large language models. The input is a student
competency profile built from the history of solved problems and errors; the output is
a structured note with text and visual materials adapted to the individual student's
level. A dedicated pipeline handles the automatic generation of mathematical
illustrations --- function graphs, geometric diagrams, and flowcharts --- ensuring
their geometric correctness.

\medskip
\noindent\textbf{Keywords:} personalized learning, educational content generation,
large language models, retrieval-augmented generation